\chapter*{Front Matter}
\addcontentsline{toc}{chapter}{Front Matter}

% Source: topping, PDF page 1
\begin{center}
\vspace*{0.33\textheight}

{\Large LECTURES ON THE RICCI FLOW\textsuperscript{1}}

\vspace{2.5em}

Peter Topping

\vspace{2.0em}

March 9, 2006

\vfill
\end{center}

\footnotetext[1]{\copyright Peter Topping 2004, 2005, 2006.}

% Source: topping, PDF page 2
Contents

\begingroup
\sloppy

1 Introduction ............................................................ 6
1.1 Ricci flow: what is it, and from where did it come? ................. 6
1.2 Examples and special solutions ....................................... 8
1.2.1 Einstein manifolds ................................................. 8
1.2.2 Ricci solitons ..................................................... 8
1.2.3 Parabolic rescaling of Ricci flows ............................... 11
1.3 Getting a feel for Ricci flow ....................................... 12
1.3.1 Two dimensions .................................................... 12
1.3.2 Three dimensions .................................................. 13
1.4 The topology and geometry of manifolds in low dimensions ............ 17
1.5 Using Ricci flow to prove topological and geometric results ......... 21

2 Riemannian geometry background ........................................ 24
2.1 Notation and conventions ............................................ 24
2.2 Einstein metrics .................................................... 28
2.3 Deformation of geometric quantities as the Riemannian metric is deformed 28
2.3.1 The formulae ...................................................... 28
2.3.2 The calculations .................................................. 32
2.4 Laplacian of the curvature tensor ................................... 39
2.5 Evolution of curvature and geometric quantities under Ricci flow .... 41

3 The maximum principle ................................................. 44
3.1 Statement of the maximum principle .................................. 44
3.2 Basic control on the evolution of curvature ......................... 45
3.3 Global curvature derivative estimates ............................... 49

4 Comments on existence theory for parabolic PDE ........................ 53
4.1 Linear scalar PDE ................................................... 53
4.2 The principal symbol ................................................ 54
4.3 Generalisation to Vector Bundles .................................... 56
4.4 Properties of parabolic equations .................................. 58

% Source: topping, PDF page 3
\section*{5 \quad Existence theory for the Ricci flow \hfill 59}
\subsection*{5.1 \quad Ricci flow is not parabolic \hfill 59}
\subsection*{5.2 \quad Short-time existence and uniqueness: The DeTurck trick \hfill 60}
\subsection*{5.3 \quad Curvature blow-up at finite-time singularities \hfill 63}

\section*{6 \quad Ricci flow as a gradient flow \hfill 67}
\subsection*{6.1 \quad Gradient of total scalar curvature and related functionals \hfill 67}
\subsection*{6.2 \quad The $\mathcal{F}$-functional \hfill 68}
\subsection*{6.3 \quad The heat operator and its conjugate \hfill 70}
\subsection*{6.4 \quad A gradient flow formulation \hfill 70}
\subsection*{6.5 \quad The classical entropy \hfill 74}
\subsection*{6.6 \quad The zeroth eigenvalue of $-4\Delta + R$ \hfill 76}

\section*{7 \quad Compactness of Riemannian manifolds and flows \hfill 78}
\subsection*{7.1 \quad Convergence and compactness of manifolds \hfill 79}
\subsection*{7.2 \quad Convergence and compactness of flows \hfill 82}
\subsection*{7.3 \quad Blowing up at singularities I \hfill 83}

\section*{8 \quad Perelman's $\mathcal{W}$ entropy functional \hfill 85}
\subsection*{8.1 \quad Definition, motivation and basic properties \hfill 85}
\subsection*{8.2 \quad Monotonicity of $\mathcal{W}$ \hfill 91}
\subsection*{8.3 \quad No local volume collapse where curvature is controlled \hfill 94}
\subsection*{8.4 \quad Volume ratio bounds imply injectivity radius bounds \hfill 100}
\subsection*{8.5 \quad Blowing up at singularities II \hfill 102}

\section*{9 \quad Curvature pinching and preserved curvature properties under Ricci flow \hfill 104}
\subsection*{9.1 \quad Overview \hfill 104}
\subsection*{9.2 \quad The Einstein Tensor, $E$ \hfill 105}
\subsection*{9.3 \quad Evolution of $E$ under the Ricci flow \hfill 106}
\subsection*{9.4 \quad The Uhlenbeck Trick \hfill 107}
\subsection*{9.5 \quad Formulae for parallel functions on vector bundles \hfill 109}
\subsection*{9.6 \quad An ODE-PDE theorem \hfill 112}
\subsection*{9.7 \quad Applications of the ODE-PDE theorem \hfill 115}

\section*{10 \quad Three-manifolds with positive Ricci curvature, and beyond. \hfill 123}
\subsection*{10.1 \quad Hamilton's theorem \hfill 123}
\subsection*{10.2 \quad Beyond the case of positive Ricci curvature \hfill 125}

\section*{A \quad Connected sum \hfill 127}

\section*{Index \hfill 132}
\endgroup

% Source: topping, PDF page 4
\begin{center}
\vspace*{0.16\textheight}

{\Huge Preface}
\end{center}

These notes represent an updated version of a course on Hamilton’s Ricci
flow that I gave at the University of Warwick in the spring of 2004. I
have aimed to give an introduction to the main ideas of the subject, a large
proportion of which are due to Hamilton over the period since he introduced
the Ricci flow in 1982. The main difference between these notes and others
which are available at the time of writing is that I follow the quite different
route which is natural in the light of work of Perelman from 2002. It is now
understood how to ‘blow up’ general Ricci flows near their singularities,
as one is used to doing in other contexts within geometric analysis. This
technique is now central to the subject, and we emphasise it throughout,
illustrating it in Chapter 10 by giving a modern proof of Hamilton’s theorem
that a closed three-dimensional manifold which carries a metric of positive
Ricci curvature is a spherical space form.

Aside from the selection of material, there is nothing in these notes which
should be considered new. There are quite a few points which have been
clarified, and we have given some proofs of well-known facts for which we
know of no good reference. The proof we give of Hamilton’s theorem does
not appear elsewhere in print, but should be clear to experts. The reader will
also find some mild reformulations, for example of the curvature pinching
results in Chapter 9.

The original lectures were delivered to a mixture of graduate students, post-
docs, staff, and even some undergraduates. Generally I assumed that the
audience had just completed a first course in differential geometry, and an
elementary course in PDE, and were just about to embark on a more ad-
vanced course in PDE. I tried to make the lectures accessible to the general
mathematician motivated by the applications of the theory to the Poincaré
conjecture, and Thurston’s geometrization conjecture (which are discussed

% Source: topping, PDF page 5
briefly in Sections 1.4 and 1.5). This has obviously affected my choice of
emphasis. I have suppressed some of the analytical issues, as discussed below,
but compiled a list of relevant Riemannian geometry calculations in Chapter 2.

There are some extremely important aspects of the theory which do not get a
mention in these notes. For example, Perelman's $\mathcal{L}$-length, which is a key
tool when developing the theory further, and Hamilton's Harnack estimates.
There is no discussion of the K\"ahler-Ricci flow.

We have stopped just short of proving the Hamilton-Ivey pinching result
which makes the study of singularities in three-dimensions tractable, although
we have covered the necessary techniques to deal with this, and may add an
exposition at a later date.

The notes are not completely self-contained. In particular, we state/use the
following without giving full proofs:

\begin{enumerate}
\item[(i)] Existence and uniqueness theory for quasilinear parabolic equations
on vector bundles. This is a long story involving rather different techniques to
those we focus on in this work. Unfortunately, it is not feasible just to quote
theorems from existing sources, and one must learn this theory for oneself;

\item[(ii)] Compactness theorems for manifolds and flows. The full proofs of
these are long, but a treatment of Ricci flow without using them would be very
misleading;

\item[(iii)] Parts of Lemma 8.1.8 which involves analysis beyond the level we were
assuming. We have given a reference, and intend to give a simple proof in later
notes.
\end{enumerate}

An updated version of these notes should be published in the L.M.S. Lecture
notes series, in conjunction with Cambridge University Press, and are also
available at:

\begin{center}
\texttt{http://www.maths.warwick.ac.uk/\string~topping/RFnotes.html}
\end{center}

Readers are invited to send comments and corrections to:

% Source: topping, PDF page 6
I would like to thank the audience of the course for making some useful
comments, especially Young Choi and Mario Micallef. Thanks also to John
Lott for comments on and typographical corrections to a 2005 version of
the notes. Parts of the original course benefited from conversations with a
number of people, including Klaus Ecker and Miles Simon. Brendan Owens
and Gero Friesecke have kindly pointed out some typographical mistakes.
Parts of the notes have been prepared whilst visiting the University of Nice,
the Albert Einstein Max-Planck Institute in Golm and Free University in
Berlin, and I would like to thank these institutions for their hospitality.
Finally, I would like to thank Neil Cource for preparing all the figures,
turning a big chunk of the original course notes into \LaTeX, and making
some corrections.
